% Part: normal-modal-logic
% Chapter: syntax-and-semantics
% Section: modal-validity

\documentclass[../../../include/open-logic-section]{subfiles}

\begin{document}

\olfileid[hi]{nml}{syn}{val}

\olsection{मान्यता}

\begin{explain}
  सभी प्रतिरूपों में, अर्थात् प्रत्येक प्रतिरूप के प्रत्येक जगत पर, सत्य
  !!{formula} विशेष रूप से रोचक हैं। वे उन मोडल प्रतिज्ञप्तियों को निरूपित करते
  हैं जो $\Box$ और $\Diamond$ की व्याख्या चाहे जैसे की जाए, सत्य रहती हैं,
  बशर्ते व्याख्या इस अर्थ में ``सामान्य'' हो कि वह संभव जगतों पर किसी
  अभिगम्यता संबंध से उत्पन्न होती हो। ऐसे !!{formula}ों को हम \emph{मान्य}
  कहते हैं। उदाहरणार्थ, $\Box(p \land q) \lif \Box p$ मान्य है। किंतु
  $\Box$ की एलेथिक व्याख्या के आधार पर जिन कुछ !!{formula}ों के मान्य होने की
  आशा की जा सकती है, जैसे $\Box p \lif p$, वे मान्य नहीं हैं। संबंधात्मक
  प्रतिरूपों की रुचिकर विशेषता यह है कि $\Box$ और $\Diamond$ की भिन्न
  व्याख्याओं को भिन्न प्रकार के अभिगम्यता संबंधों से ग्रहण किया जा सकता है।
  इससे संकेत मिलता है कि मान्यता को केवल \emph{सभी} प्रतिरूपों के सापेक्ष ही
  नहीं, बल्कि \emph{किसी निश्चित प्रकार के} सभी प्रतिरूपों के सापेक्ष भी
  परिभाषित करना चाहिए। उदाहरणार्थ, यह निकलेगा कि $\Box p \lif p$ उन सभी
  प्रतिरूपों में सत्य है जिनमें प्रत्येक जगत स्वयं से अभिगम्य है, अर्थात्
  $R$ स्वपरावर्ती है। प्रतिरूपों के वर्गों के सापेक्ष मान्यता की परिभाषा
  हमें इसे संक्षेप में कहने देती है: $\Box p \lif p$ स्वपरावर्ती प्रतिरूपों
  के वर्ग में मान्य है।
\end{explain}

\begin{defn}
  कोई !!{formula}~$!A$, प्रतिरूपों के वर्ग $\mClass{C}$ में \emph{मान्य} है,
  यदि वह $\mClass{C}$ के प्रत्येक प्रतिरूप में (अर्थात् प्रत्येक प्रतिरूप के
  प्रत्येक जगत पर) सत्य हो। यदि $!A$, $\mClass{C}$ में मान्य है, तो हम
  $\mClass{C} \Entails !A$ लिखते हैं; और यदि $!A$, \emph{सभी} प्रतिरूपों
  के वर्ग में मान्य है, तो $\Entails !A$ लिखते हैं।
\end{defn}

\begin{prop}\ollabel{prop:subset-class}
  यदि $!A$, $\mClass{C}$ में मान्य है, तो वह प्रत्येक वर्ग
  $\mClass{C}' \subseteq \mClass{C}$ में भी मान्य है।
\end{prop}

\begin{prop}\ollabel{prop:Nec-rule}
  यदि $!A$ मान्य है, तो $\Box!A$ भी मान्य है।
\end{prop}

\begin{proof}
  मान लें $\Entails !A$। $\Entails \Box!A$ दिखाने के लिए कोई प्रतिरूप
  $\mModel{M} = \tuple{W, R, V}$ और $w \in W$ लें। यदि $Rww'$, तो $!A$ की
  मान्यता से $\mSat{M}{!A}[w']$, अतः $\mSat{M}{\Box!A}[w]$ भी। चूँकि
  $\mModel{M}$ और $w$ स्वेच्छ थे, $\Entails \Box!A$।
\end{proof}

\begin{prob}
  दिखाइए कि निम्नलिखित मान्य हैं:
  \begin{enumerate}
  \item $\Box p \lif \Box (q \lif p)$;
  \item $\Box \lnot \lfalse$;
  \item $\Box p \lif (\Box q \lif \Box p)$।
  \end{enumerate}
\end{prob}

\begin{prob}
  दिखाइए कि $!A \lif \Box!A$, उन प्रतिरूपों
  $\mModel{M} = \tuple{W, R, V}$ के वर्ग $\mClass{C}$ में मान्य है जिनमें
  $W = \{w\}$। इसी तरह दिखाइए कि $!B \lif \Box !A$ और
  $\Diamond !A \lif !B$, उन प्रतिरूपों $\mModel{M} = \tuple{W, R, V}$ के
  वर्ग में मान्य हैं जिनमें $R = \emptyset$।
\end{prob}

\end{document}
